\documentclass[11pt,a4paper]{article}

\usepackage[utf8]{inputenc}
\usepackage[margin=0.9in]{geometry}
\usepackage{amsmath,amssymb,amsthm,amsfonts}
\usepackage{mathtools}
\usepackage{hyperref}
\usepackage{booktabs}
\usepackage{cite}
\usepackage{microtype}
\usepackage{array}

\hypersetup{
    colorlinks=true,
    linkcolor=blue,
    citecolor=red,
    urlcolor=blue
}

% Theorem Environments
\theoremstyle{plain}
\newtheorem{theorem}{Theorem}[section]
\newtheorem{lemma}[theorem]{Lemma}
\newtheorem{proposition}[theorem]{Proposition}
\newtheorem{corollary}[theorem]{Corollary}
\newtheorem{conjecture}[theorem]{Conjecture}

\theoremstyle{definition}
\newtheorem{definition}[theorem]{Definition}
\newtheorem{example}[theorem]{Example}

\theoremstyle{remark}
\newtheorem{remark}[theorem]{Remark}

\numberwithin{equation}{section}

\title{\textbf{Nonlinear Simplicial Waves and Anomalous Porous Transport Induced by Beta-Kernel Fractional Laplacians}}
\author{\textbf{Reinaldo Maia Silva-Filho} \\ \textit{Universidade Federal de Lavras (UFLA)}}
\date{\today}


\DeclareMathOperator*{\argmin}{argmin}
\providecommand{\doi}[1]{\href{https://doi.org/#1}{\texttt{doi:#1}}}
\providecommand{\Hyp}{\mathbb{H}}
\providecommand{\Sph}{\mathbb{S}}
\providecommand{\II}{\mathrm{II}}
\providecommand{\R}{\mathbb{R}}
\providecommand{\Area}{\mathrm{Area}}
\providecommand{\Hsn}{\mathcal{H}}
\providecommand{\BV}{\mathrm{BV}}
\providecommand{\TV}{\mathrm{TV}}
\providecommand{\cutnorm}[1]{\|#1\|_{\square}}

\begin{document}

\maketitle

\begin{abstract}
We develop a comprehensive analytical and mathematical physics framework for non-local wave dynamics and subdiffusive transport governed by the \emph{Fractional Simplicial Laplacian} $\Delta_{\Delta_m}^\alpha$, an operator induced by continuous multinomial Beta-kernels on the standard simplex $\Delta_{m-1}(\alpha)$. In the first part of this work, we formulate the \emph{Nonlinear Simplicial Schr\"odinger Equation} (NLSE). We prove the global conservation of particle mass $\mathcal{N}[\psi]$ and simplicial Hamiltonian energy $\mathcal{E}[\psi]$ via the self-adjointness of the continuous Beta-integral kernel. A modulational instability analysis shows that, because the symbol is bounded, the unstable band can cover all non-zero wavevectors, unlike in the local NLSE. Also because the symbol is bounded, the energy at fixed mass is unbounded below in the focusing case, so there are no energy-minimizing ground states; the existence and symmetry of solitary waves are stated as a conjecture. In the second part, we formulate the space-time fractional simplicial diffusion equation for anomalous transport in anisotropic media. Using Fourier--Laplace methods, we give the exact propagator in terms of Mittag-Leffler functions and prove that the mean squared displacement tensor is exactly $\langle \mathbf{x}\mathbf{x}^T \rangle(t) = \frac{\mathcal{K}_{\text{diff}}}{\alpha^2 \Gamma(\beta+1)} \mathbf{M}_2(\alpha)\, t^\beta$, where $\mathbf{M}_2(\alpha) = a\,\mathrm{Id} + b\,\mathbf{1}\mathbf{1}^T$ is the second-moment matrix of the jump kernel.
\end{abstract}

\tableofcontents
\newpage

% ==============================================================================
\section{Introduction and Physical Motivation}
% ==============================================================================

Non-local partial differential equations (PDEs) and fractional calculus have become foundational tools for describing memory-dependent dynamics, anomalous transport, and non-local wave phenomena across quantum physics, fluid mechanics, and geophysical transport \cite{podlubny1998, samko1993, metzler2000, laskin2000}. In traditional fractional quantum mechanics and anomalous diffusion models, the non-local spatial operator is typically chosen as the isotropic Riesz fractional Laplacian $(-\Delta)^{\alpha/2}$, whose integral representation:
\begin{equation}
(-\Delta)^{\alpha/2} u(\mathbf{x}) = C_{d, \alpha} \operatorname{P.V.} \int_{\mathbb{R}^d} \frac{u(\mathbf{x}) - u(\mathbf{y})}{|\mathbf{x} - \mathbf{y}|^{d+\alpha}} \, d\mathbf{y},
\end{equation}
is isotropic, radially symmetric under the orthogonal group $O(d)$, and characterized by algebraic power-law tails that extend across the entire space $\mathbb{R}^d$.

However, in many realistic physical media—such as anisotropic fractured rocks, crystalline waveguides, packed colloidal arrays, and simplicial metamaterials—the underlying microstructure imposes two fundamental physical constraints:
\begin{enumerate}
    \item \textbf{Directional Geometric Anisotropy:} Transport occurs predominantly along structured interconnected channels with characteristic intersection angles (such as triangular or tetrahedral lattices), breaking continuous rotational symmetry.
    \item \textbf{Finite Horizon of Non-Local Steps:} Microscopic jump lengths per collision event are geometrically bounded by the pore or cell size, precluding unphysical infinite-velocity jump tails.
\end{enumerate}

To bridge this fundamental gap, we introduce and analyze dynamical and transport processes governed by the \emph{Fractional Simplicial Laplacian} $\Delta_{\Delta_m}^\alpha$, constructed from the continuous multinomial Beta-kernel on the standard $(m-1)$-simplex $\Delta_{m-1}(\alpha)$.

% ==============================================================================
\section{The Fractional Simplicial Laplacian and its Metric Geometry}
% ==============================================================================

Let $\mathbf{x} = (x_1, \dots, x_{m-1}) \in \mathbb{R}^{m-1}$ denote spatial coordinates, and let $\Delta_{m-1}(\alpha) \subset \mathbb{R}_{\ge 0}^{m-1}$ be the standard scaled simplex of size $\alpha > 0$:
\begin{equation}
\Delta_{m-1}(\alpha) \coloneqq \left\{\mathbf{y} \in \mathbb{R}_{\ge 0}^{m-1} \;\middle|\; \sum_{j=1}^{m-1} y_j \le \alpha\right\}.
\end{equation}
The dependent barycentric coordinate is defined by $y_m \coloneqq \alpha - \sum_{j=1}^{m-1} y_j$. The continuous multinomial Beta-kernel is defined by
\begin{equation}
\binom{\alpha}{\mathbf{y}} \coloneqq \frac{\Gamma(\alpha+1)}{\prod_{j=1}^m \Gamma(y_j+1)}, \quad I_m(\alpha) \coloneqq \int_{\Delta_{m-1}(\alpha)} \binom{\alpha}{\mathbf{y}} \, d\mathbf{y}.
\end{equation}

\begin{definition}[Fractional Simplicial Laplacian]
For fractional parameter $\alpha > 0$, the Fractional Simplicial Laplacian acting on $u \in L^2(\mathbb{R}^{m-1})$ is defined by
\begin{equation}
\label{eq:simplicial_laplacian_def}
\Delta_{\Delta_m}^\alpha u(\mathbf{x}) \coloneqq \frac{1}{2\alpha^2 I_m(\alpha)}\int_{\Delta_{m-1}(\alpha)} \binom{\alpha}{\mathbf{y}} \Big[u(\mathbf{x}+\mathbf{y}) - 2u(\mathbf{x}) + u(\mathbf{x}-\mathbf{y})\Big] d\mathbf{y}.
\end{equation}
\end{definition}

\begin{remark}[Identification of Fractional Order and Support Radius]
As introduced in the canonical formalism \cite{silvafilho2026simplex}, the parameter $\alpha$ is employed dually as the fractional order of the continuous multinomial coefficient and as the geometric support radius defining the integration domain $\Delta_{m-1}(\alpha)$. This explicitly couples the non-local diffusion strength to its macroscopic geometric extent as a phenomenological ansatz. Generalizing this operator to decouple the combinatorial fractional index from an independent geometric scale parameter remains an open area for future formalization.

\emph{Dimensional convention.} Since $\alpha$ enters $\Gamma(\alpha+1)$, it is dimensionless. Its use as a support radius therefore presupposes a fixed unit of length $\ell_0$: physical positions are $\ell_0\mathbf{x}$, the jump kernel has support radius of order $\alpha\ell_0$, and the symbol $\sigma^\alpha_{\Delta_m}$ is measured in units of $\ell_0^{-2}$. With this convention, $\frac{\hbar^2}{2M}\sigma^\alpha_{\Delta_m}$ is an energy, and \eqref{eq:nlse_main} below is dimensionally consistent. The choice of $\ell_0$ is part of the model and is not fixed by the kernel.
\end{remark}

\begin{proposition}[Self-Adjointness and Positivity]
The operator $-\Delta_{\Delta_m}^\alpha$ is linear, symmetric, and positive semi-definite on $L^2(\mathbb{R}^{m-1})$:
\begin{equation}
\langle u, -\Delta_{\Delta_m}^\alpha u \rangle_{L^2} = \frac{1}{2\alpha^2 I_m(\alpha)}\int_{\mathbb{R}^{m-1}} d\mathbf{x} \int_{\Delta_{m-1}(\alpha)} \binom{\alpha}{\mathbf{y}} |u(\mathbf{x}+\mathbf{y}) - u(\mathbf{x})|^2 d\mathbf{y} \ge 0.
\end{equation}
\end{proposition}

\begin{proof}
Expanding the inner product $\langle u, -\Delta_{\Delta_m}^\alpha u \rangle_{L^2} = \int_{\mathbb{R}^{m-1}} u^*(\mathbf{x}) (-\Delta_{\Delta_m}^\alpha u(\mathbf{x})) d\mathbf{x}$:
\begin{align*}
\langle u, -\Delta_{\Delta_m}^\alpha u \rangle &= \frac{1}{2\alpha^2 I_m(\alpha)}\int_{\Delta_{m-1}(\alpha)} \binom{\alpha}{\mathbf{y}} \notag \\
&\quad \times \left[ \int_{\mathbb{R}^{m-1}} \Big(2|u(\mathbf{x})|^2 - u^*(\mathbf{x})u(\mathbf{x}+\mathbf{y}) - u^*(\mathbf{x})u(\mathbf{x}-\mathbf{y})\Big) d\mathbf{x} \right] d\mathbf{y}.
\end{align*}
Using the translation invariance $\int u^*(\mathbf{x})u(\mathbf{x}-\mathbf{y})d\mathbf{x} = \int u^*(\mathbf{x}+\mathbf{y})u(\mathbf{x})d\mathbf{x}$:
\begin{equation}
\int_{\mathbb{R}^{m-1}} \Big(2|u(\mathbf{x})|^2 - u^*(\mathbf{x})u(\mathbf{x}+\mathbf{y}) - u(\mathbf{x})u^*(\mathbf{x}+\mathbf{y})\Big) d\mathbf{x} = \int_{\mathbb{R}^{m-1}} |u(\mathbf{x}+\mathbf{y}) - u(\mathbf{x})|^2 d\mathbf{x}.
\end{equation}
Since $\binom{\alpha}{\mathbf{y}} > 0$ on the interior of $\Delta_{m-1}(\alpha)$, the integral is non-negative. It vanishes if and only if $u = 0$ in $L^2$. Indeed, vanishing forces $\int|u(\mathbf{x}+\mathbf{y}) - u(\mathbf{x})|^2d\mathbf{x} = 0$ for almost every $\mathbf{y}$ in the open set $\operatorname{int}\Delta_{m-1}(\alpha)$. By Plancherel, $\widehat{u}(\mathbf{k})(e^{i\mathbf{k}\cdot\mathbf{y}} - 1) = 0$ for almost every $\mathbf{k}$ and all such $\mathbf{y}$. For $\mathbf{k} \ne \mathbf{0}$, $e^{i\mathbf{k}\cdot\mathbf{y}} \ne 1$ for some $\mathbf{y}$ in the open set, so $\widehat{u} = 0$ almost everywhere. (Non-zero constants are not in $L^2(\mathbb{R}^{m-1})$.)
\end{proof}

\subsection{Boundedness and Relation to Spectral Fractional Laplacians}

\begin{remark}[Boundedness]
Because $\binom{\alpha}{\mathbf{y}}$ is integrable with compact support, $-\Delta_{\Delta_m}^\alpha$ is a bounded operator on $L^2(\mathbb{R}^{m-1})$ with $0 \le -\Delta_{\Delta_m}^\alpha \le 2/\alpha^2$ \cite{silvafilho2026simplex}. No boundary conditions or domain restrictions are needed to make it self-adjoint.
\end{remark}

\begin{remark}[Spectral fractional Laplacians on bounded domains are a different object]
\label{rem:spectral_fractional}
One may also define, on a bounded domain $\Omega$ (for instance a simplex), the \emph{spectral} fractional Laplacian $(-\Delta_D)^s u = \sum_n \lambda_n^s \langle u, \phi_n\rangle \phi_n$ from the Dirichlet eigenpairs $(\lambda_n,\phi_n)$ of $-\Delta$. Two facts are worth recording.
(a) For $s \in (0,1)$, the spectral operator does \emph{not} coincide with the restricted singular-integral fractional Laplacian on $\Omega$; the two differ on bounded domains, and even their eigenvalues differ \cite{servadei2014spectrum}.
(b) For $s = 2$, the domain of $(-\Delta_D)^2$ is $\{u \in \operatorname{Dom}(-\Delta_D) : \Delta u \in \operatorname{Dom}(-\Delta_D)\}$. For a bounded convex domain, and in particular for a simplex, $\operatorname{Dom}(-\Delta_D) = H^2(\Omega) \cap H^1_0(\Omega)$ \cite{grisvard1985elliptic}, so the domain is $\{u \in H^2 \cap H^1_0 : \Delta u \in H^2 \cap H^1_0\}$. This encodes the Navier conditions $u = \Delta u = 0$ on $\partial\Omega$, not the clamped conditions $u = \partial_{\mathbf{n}} u = 0$, which define a different fourth-order operator \cite{gazzola2010polyharmonic}. For smooth $\partial\Omega$ the domain equals $\{u \in H^4(\Omega) : u = \Delta u = 0 \text{ on } \partial\Omega\}$. On polygons this depends on the angles. On a rectangle or an equilateral triangle, where odd reflections across the sides tile the plane, the same equality holds. On a general triangle it fails: near a corner of angle $\omega$ with $\pi/\omega < 3$ not an integer, the Dirichlet problem has the singular solution $r^{\pi/\omega}\sin(\pi\theta/\omega)$, which is not in $H^4$ \cite{grisvard1985elliptic}. Thus on a general simplex $H^4$ regularity fails near edges and vertices.
Neither operator is the Beta-kernel operator $\Delta_{\Delta_m}^\alpha$ studied in this chapter.
\end{remark}

\begin{theorem}[Dispersion Symbol and Long-Wavelength Metric; {\cite[Thm.~7.2]{silvafilho2026simplex}}]
\label{thm:dispersion}
The Fourier symbol defined by $-\Delta_{\Delta_m}^\alpha e^{i\mathbf{k}\cdot\mathbf{x}} = \sigma_{\Delta_m}^\alpha(\mathbf{k}) e^{i\mathbf{k}\cdot\mathbf{x}}$ is
\begin{equation}
\sigma_{\Delta_m}^\alpha(\mathbf{k}) = \frac{1}{\alpha^2 I_m(\alpha)}\int_{\Delta_{m-1}(\alpha)} \binom{\alpha}{\mathbf{y}}\big(1 - \cos(\mathbf{k}\cdot\mathbf{y})\big)\,d\mathbf{y} \in \Big[0, \frac{2}{\alpha^2}\Big],
\end{equation}
an even entire function of $\mathbf{k}$ (the Fourier transform of a compactly supported density). As $|\mathbf{k}| \to 0$,
\begin{equation}
\sigma_{\Delta_m}^\alpha(\mathbf{k}) = \frac{1}{2\alpha^2}\mathbf{k}^T \mathbf{M}_2(\alpha)\mathbf{k} + \mathcal{O}(|\mathbf{k}|^4), \qquad \mathbf{M}_2(\alpha) = a(\alpha)\,\mathrm{Id}_{m-1} + b(\alpha)\,\mathbf{1}\mathbf{1}^T,
\end{equation}
where $\mathbf{M}_2$ is the second-moment matrix of the normalized kernel. (The closed form $\frac{1}{\alpha^2}[1 - R^\alpha\cos(\alpha\Theta)]$ belongs to the lattice multinomial, not to this continuous kernel; see \cite{silvafilho2026simplex}. Even for the lattice multinomial the small-$\mathbf{k}$ limit is $\frac{S^2}{2m^2} + \frac{Q - S^2/m}{2m\alpha}$, with $S = \sum_j k_j$ and $Q = |\mathbf{k}|^2$. It reduces to $\frac{1}{2m\alpha}\mathbf{k}^T\mathbf{A}_{m-1}\mathbf{k}$, with $\mathbf{A}_{m-1} = \mathrm{Id} +\mathbf{1}\mathbf{1}^T$ the Gram matrix of Chapter~3, only on the constraint $\sum_j k_j = 0$. Both forms then equal $\frac{Q}{2m\alpha}$.)
\end{theorem}

% ==============================================================================
\section{Nonlinear Simplicial Schr\"odinger Equations (NLSE)}
% ==============================================================================

We formulate the Cauchy problem for the non-local Fractional Simplicial NLSE on $\mathbb{R}^{m-1} \times \mathbb{R}^+$:
\begin{equation}
\label{eq:nlse_main}
i \hbar \frac{\partial \psi(\mathbf{x}, t)}{\partial t} = -\frac{\hbar^2}{2M}\Delta_{\Delta_m}^\alpha \psi(\mathbf{x}, t) + V(\mathbf{x})\psi(\mathbf{x}, t) - \kappa |\psi(\mathbf{x}, t)|^{2p}\psi(\mathbf{x}, t),
\end{equation}
where $M > 0$ is the effective mass, $V(\mathbf{x})$ is an external real-valued potential, $\kappa \in \mathbb{R}$ is the non-linear coupling constant ($\kappa > 0$ focusing, $\kappa < 0$ defocusing), and $p > 0$ is the non-linearity exponent.

\subsection{Fundamental Conservation Laws}

\begin{theorem}[Global Conservation of Mass and Energy]
Let $V$ be real-valued and independent of $t$, and let $\psi(\mathbf{x}, t)$ be a smooth solution to \eqref{eq:nlse_main} decaying sufficiently fast at spatial infinity. Use the unitary Fourier transform $\widehat{\psi}(\mathbf{k}) = (2\pi)^{-(m-1)/2}\int \psi(\mathbf{x})e^{-i\mathbf{k}\cdot\mathbf{x}}d\mathbf{x}$, so that $\int \sigma |\widehat\psi|^2 d\mathbf{k} = \langle \psi, -\Delta^\alpha_{\Delta_m}\psi\rangle$. Then:
\begin{enumerate}
    \item \textbf{Conservation of Total Mass (Particle Number):}
    \begin{equation}
    \mathcal{N}[\psi(t)] \coloneqq \int_{\mathbb{R}^{m-1}} |\psi(\mathbf{x}, t)|^2 \, d\mathbf{x} = \mathcal{N}[\psi(0)].
    \end{equation}
    \item \textbf{Conservation of Simplicial Energy (Hamiltonian):}
    \begin{align}
    \mathcal{E}[\psi(t)] &\coloneqq \frac{\hbar^2}{2M}\int_{\mathbb{R}^{m-1}} \sigma_{\Delta_m}^\alpha(\mathbf{k}) |\widehat{\psi}(\mathbf{k}, t)|^2 d\mathbf{k} + \int_{\mathbb{R}^{m-1}} V(\mathbf{x})|\psi(\mathbf{x}, t)|^2 d\mathbf{x} \nonumber\\
    &\quad - \frac{\kappa}{p+1}\int_{\mathbb{R}^{m-1}} |\psi(\mathbf{x}, t)|^{2p+2} d\mathbf{x} = \mathcal{E}[\psi(0)].
    \end{align}
\end{enumerate}
\end{theorem}

\begin{proof}
\leavevmode
\begin{enumerate}
    \item Multiplying \eqref{eq:nlse_main} by $\psi^*$ and taking the imaginary part:
    \begin{equation}
    \frac{\hbar}{2}\frac{\partial |\psi|^2}{\partial t} = \operatorname{Im}\left(-\frac{\hbar^2}{2M}\psi^* \Delta_{\Delta_m}^\alpha \psi\right).
    \end{equation}
    Integrating over $\mathbb{R}^{m-1}$ and applying the self-adjointness of $\Delta_{\Delta_m}^\alpha$:
    \begin{equation}
    \frac{d}{dt}\mathcal{N}[\psi(t)] = -\frac{\hbar}{M}\operatorname{Im}\langle \psi, \Delta_{\Delta_m}^\alpha \psi \rangle_{L^2} = 0,
    \end{equation}
    since $\langle \psi, \Delta_{\Delta_m}^\alpha \psi \rangle_{L^2} \in \mathbb{R}$.
    
    \item Differentiating $\mathcal{E}[\psi(t)]$ with respect to $t$:
    \begin{equation}
    \frac{d\mathcal{E}}{dt} = 2\operatorname{Re}\int_{\mathbb{R}^{m-1}} \psi_t^* \left[ -\frac{\hbar^2}{2M}\Delta_{\Delta_m}^\alpha \psi + V(\mathbf{x})\psi - \kappa |\psi|^{2p}\psi \right] d\mathbf{x}.
    \end{equation}
    Substituting the equation of motion $i\hbar \psi_t = \mathcal{H}\psi$:
    \begin{equation}
    \frac{d\mathcal{E}}{dt} = 2\operatorname{Re}\int_{\mathbb{R}^{m-1}} \left(\frac{1}{i\hbar}\mathcal{H}\psi\right)^* (\mathcal{H}\psi) \, d\mathbf{x} = 2\operatorname{Re}\left(-\frac{1}{i\hbar}\|\mathcal{H}\psi\|_{L^2}^2\right) = 0.
    \end{equation}
\end{enumerate}
\end{proof}

\subsection{Modulational Instability}

Let $V \equiv 0$. Consider a uniform plane-wave background state with amplitude $\psi_0(\mathbf{x}, t) = \sqrt{\rho_0} e^{-i \mu_0 t / \hbar}$, where the chemical potential is $\mu_0 = -\kappa \rho_0^p$. We perturb the background as
\begin{equation}
\psi(\mathbf{x}, t) = \left(\sqrt{\rho_0} + \delta u(\mathbf{x}, t) + i \delta v(\mathbf{x}, t)\right) e^{-i \mu_0 t / \hbar}.
\end{equation}

\begin{theorem}[Simplicial Modulational Instability]
Linearizing \eqref{eq:nlse_main} with respect to perturbations $(\delta u, \delta v) \propto e^{i(\mathbf{k}\cdot\mathbf{x} - \Omega t)}$, the perturbation frequency satisfies the dispersion relation:
\begin{equation}
\hbar^2 \Omega^2(\mathbf{k}) = \frac{\hbar^2 \sigma_{\Delta_m}^\alpha(\mathbf{k})}{2M} \left[ \frac{\hbar^2 \sigma_{\Delta_m}^\alpha(\mathbf{k})}{2M} - 2\kappa p \rho_0^p \right].
\end{equation}
In the focusing regime ($\kappa > 0$), the plane wave is modulationally unstable ($\Omega^2 < 0$) for wavevectors satisfying
\begin{equation}
\sigma_{\Delta_m}^\alpha(\mathbf{k}) < \frac{4M\kappa p \rho_0^p}{\hbar^2},
\end{equation}
with maximum instability growth rate $\gamma_{\max} = \frac{\kappa p \rho_0^p}{\hbar}$ attained on the resonant set $\{\sigma_{\Delta_m}^\alpha(\mathbf{k}) = 2M\kappa p\rho_0^p/\hbar^2\}$. This set is non-empty when $2M\kappa p\rho_0^p/\hbar^2 < \sup_{\mathbf{k}}\sigma_{\Delta_m}^\alpha$ ($\le 2/\alpha^2$), by continuity of $\sigma$ and $\sigma(\mathbf{0}) = 0$. In the case of equality it is non-empty exactly when the supremum is attained. Since the symbol is bounded, for $4M\kappa p\rho_0^p/\hbar^2 > \sup\sigma$ \emph{every} wavevector $\mathbf{k} \ne \mathbf{0}$ is unstable, unlike the local NLSE ($\mathbf{k} = \mathbf{0}$ gives $\Omega = 0$).
\end{theorem}
\begin{proof}
Linearizing gives $\hbar\,\partial_t \delta u = \varepsilon\,\delta v$ and $\hbar\,\partial_t \delta v = -(\varepsilon - 2\kappa p\rho_0^p)\,\delta u$ in Fourier variables, with $\varepsilon = \hbar^2\sigma/(2M)$. Hence $\hbar^2\Omega^2 = \varepsilon(\varepsilon - 2\kappa p\rho_0^p)$, which is minimized at $\varepsilon = \kappa p\rho_0^p$, giving $\hbar^2\Omega^2 = -(\kappa p\rho_0^p)^2$.
\end{proof}

\subsection{Simplicial Solitary Waves and Anisotropic Solitons}

For stationary solitary wave states $\psi(\mathbf{x}, t) = e^{-i\mu t/\hbar} \phi(\mathbf{x})$ with $\mu < 0$:
\begin{equation}
\label{eq:soliton_profile}
-\frac{\hbar^2}{2M}\Delta_{\Delta_m}^\alpha \phi(\mathbf{x}) + (V(\mathbf{x}) - \mu)\phi(\mathbf{x}) = \kappa |\phi(\mathbf{x})|^{2p}\phi(\mathbf{x}).
\end{equation}

\begin{conjecture}[Existence and Symmetry of Simplicial Solitary Waves]
Let $V \equiv 0$, $\kappa > 0$ and $\mu < 0$. Then \eqref{eq:soliton_profile} has a non-trivial solution $\phi \in L^2(\mathbb{R}^{m-1}) \cap L^\infty(\mathbb{R}^{m-1})$ which, after translation, is invariant under the group generated by the permutations of $x_1, \dots, x_{m-1}$ and $\mathbf{x} \mapsto -\mathbf{x}$. This group, $S_{m-1} \times \mathbb{Z}_2$, preserves $\Delta_{\Delta_m}^\alpha$. The full group $S_m$ does not act linearly: exchanging $y_m$ with $y_j$ is an affine map on $\mathbb{R}^{m-1}$, and it does not preserve the second-moment matrix $\mathbf{M}_2$ \cite[Rem.~7.3]{silvafilho2026simplex}. We do not know whether $S_{m-1} \times \mathbb{Z}_2$ is the whole group of linear symmetries of the operator. Such solutions are not radially symmetric.
\end{conjecture}

\begin{remark}
The conjecture cannot be approached by minimizing the energy at fixed mass, because in the focusing case $\kappa > 0$ that infimum is $-\infty$ for every $p > 0$. The kinetic term is bounded, $0 \le \frac{\hbar^2}{2M}\int\sigma|\widehat\psi|^2 \le \frac{\hbar^2}{M\alpha^2}\mathcal{N}[\psi]$, while for $\psi_w(\mathbf{x}) = w^{-(m-1)/2}\phi(\mathbf{x}/w)$ the mass is independent of $w$ and $\int|\psi_w|^{2p+2} = w^{-(m-1)p}\int|\phi|^{2p+2} \to \infty$ as $w \to 0$. Hence ``ground state'' must be understood as a solution of \eqref{eq:soliton_profile} selected by another principle, for instance a least-action solution on a Nehari-type set, and not as an energy minimizer. Neither existence nor symmetry is proved here. For a kernel with compact support, the decay of solitary waves is expected to be exponential rather than algebraic.
\end{remark}

% ==============================================================================
\section{Anomalous Diffusion in Anisotropic Porous Media}
% ==============================================================================

\subsection{Model Formulation}

In heterogeneous porous and fractured geological media, microscopic transport is characterized by: (i) temporal retention in pore dead-ends (subdiffusion), and (ii) anisotropic spatial hopping constrained by simplicial pore geometries.

\begin{definition}[Space-Time Fractional Simplicial Diffusion Equation]
Let $u(\mathbf{x}, t)$ denote the solute concentration. The governing equation is
\begin{equation}
\label{eq:fract_diff}
\frac{\partial^\beta u(\mathbf{x}, t)}{\partial t^\beta} = \mathcal{K}_{\text{diff}} \Delta_{\Delta_m}^\alpha u(\mathbf{x}, t), \quad \mathbf{x} \in \mathbb{R}^{m-1}, \; t > 0,
\end{equation}
with initial condition $u(\mathbf{x}, 0) = u_0(\mathbf{x})$, where $\beta \in (0, 1]$ is the temporal order of the Caputo fractional derivative:
\begin{equation}
\frac{\partial^\beta u}{\partial t^\beta}(t) \coloneqq \frac{1}{\Gamma(1-\beta)}\int_0^t (t-\tau)^{-\beta} \frac{\partial u}{\partial \tau}(\tau) \, d\tau.
\end{equation}
\end{definition}

\subsection{Exact Propagator via Mittag-Leffler Functions}

\begin{theorem}[Closed-Form Solution in Frequency Domain]
The spatial Fourier transform of the exact solution to \eqref{eq:fract_diff} is given by
\begin{equation}
\widehat{u}(\mathbf{k}, t) = E_\beta\left(-\mathcal{K}_{\text{diff}}\,\sigma_{\Delta_m}^\alpha(\mathbf{k})\,t^\beta\right) \widehat{u}_0(\mathbf{k}),
\end{equation}
where $E_\beta(z) \coloneqq \sum_{n=0}^\infty \frac{z^n}{\Gamma(\beta n + 1)}$ is the Mittag-Leffler function. Since $\sigma^\alpha_{\Delta_m}(\mathbf{k}) \to \alpha^{-2}$ as $|\mathbf{k}| \to \infty$, the fundamental solution ($u_0 = \delta$) is not a function. It is a measure with an atom $E_\beta(-\mathcal{K}_{\text{diff}}\,t^\beta/\alpha^2)\,\delta(\mathbf{x})$ at the origin, which decays only slowly in $t$ (like $t^{-\beta}$ for $\beta < 1$), plus the inverse transform of the remainder.
\end{theorem}

\begin{proof}
Applying the Laplace transform $\mathcal{L}\{f\}(s) = \int_0^\infty e^{-st}f(t)dt$ to the Caputo derivative yields $\mathcal{L}\{\partial_t^\beta u\}(s) = s^\beta \widetilde{u}(s) - s^{\beta-1}u(0)$. Applying the spatial Fourier transform $\mathbf{x} \to \mathbf{k}$:
\begin{equation}
s^\beta \widehat{\widetilde{u}}(\mathbf{k}, s) - s^{\beta-1}\widehat{u}_0(\mathbf{k}) = -\mathcal{K}_{\text{diff}}\,\sigma_{\Delta_m}^\alpha(\mathbf{k})\,\widehat{\widetilde{u}}(\mathbf{k}, s).
\end{equation}
Solving algebraically for $\widehat{\widetilde{u}}(\mathbf{k}, s)$:
\begin{equation}
\widehat{\widetilde{u}}(\mathbf{k}, s) = \frac{s^{\beta-1}}{s^\beta + \mathcal{K}_{\text{diff}}\,\sigma_{\Delta_m}^\alpha(\mathbf{k})} \widehat{u}_0(\mathbf{k}).
\end{equation}
Applying the inverse Laplace transform using the standard identity $\mathcal{L}^{-1}\left\{\frac{s^{\beta-1}}{s^\beta + a}\right\}(t) = E_\beta(-a t^\beta)$ immediately yields the result.
\end{proof}

\subsection{Mean Squared Displacement (MSD) and Simplicial Permeability Tensor}

\begin{theorem}[Anisotropic Mean Squared Displacement]
\label{thm:msd}
Assume a localized initial condition $u_0(\mathbf{x}) = \delta(\mathbf{x})$. The first and second moments of the concentration profile satisfy:
\begin{enumerate}
    \item \textbf{Mean Drift:} $\langle \mathbf{x}(t) \rangle = \mathbf{0}$.
    \item \textbf{Mean Squared Displacement Covariance Tensor:}
    \begin{equation}
    \langle \mathbf{x} \mathbf{x}^T \rangle(t) \coloneqq \int_{\mathbb{R}^{m-1}} \mathbf{x}\mathbf{x}^T u(\mathbf{x}, t)\,d\mathbf{x} = \frac{\mathcal{K}_{\text{diff}}}{\alpha^2 \, \Gamma(\beta + 1)} \mathbf{M}_2(\alpha) \cdot t^\beta,
    \end{equation}
    exactly for all $t > 0$, where $\mathbf{M}_2(\alpha) = a(\alpha)\mathrm{Id} + b(\alpha)\mathbf{1}\mathbf{1}^T$ is the second-moment matrix of Theorem~\ref{thm:dispersion}.
\end{enumerate}
\end{theorem}

\begin{proof}
Using moment-generating properties of the Fourier transform:
\begin{equation}
\langle x_j \rangle(t) = \left. i \frac{\partial \widehat{u}(\mathbf{k}, t)}{\partial k_j} \right|_{\mathbf{k}=\mathbf{0}} = \left. i E_\beta'\left(-\mathcal{K}_{\text{diff}}\sigma_{\Delta_m}^\alpha t^\beta\right) \left(-\mathcal{K}_{\text{diff}} t^\beta \frac{\partial \sigma_{\Delta_m}^\alpha}{\partial k_j}\right) \right|_{\mathbf{k}=\mathbf{0}} = 0,
\end{equation}
since $\sigma_{\Delta_m}^\alpha$ is even in $\mathbf{k}$ (the operator is built from the symmetric second difference), so $\left.\nabla_\mathbf{k}\sigma_{\Delta_m}^\alpha\right|_{\mathbf{k}=\mathbf{0}} = \mathbf{0}$.

For the second moment tensor:
\begin{equation}
\langle x_i x_j \rangle(t) = \left. -\frac{\partial^2 \widehat{u}(\mathbf{k}, t)}{\partial k_i \partial k_j} \right|_{\mathbf{k}=\mathbf{0}} = \left. -E_\beta'(0) \left(-\mathcal{K}_{\text{diff}} t^\beta \frac{\partial^2 \sigma_{\Delta_m}^\alpha}{\partial k_i \partial k_j}\right) \right|_{\mathbf{k}=\mathbf{0}}.
\end{equation}
Since $E_\beta(z) = \frac{1}{\Gamma(1)} + \frac{z}{\Gamma(\beta+1)} + \mathcal{O}(z^2)$, we have $E_\beta'(0) = \frac{1}{\Gamma(\beta+1)}$. By Theorem~\ref{thm:dispersion}, the Hessian of the symbol at $\mathbf{k}=\mathbf{0}$ is:
\begin{equation}
\left. \frac{\partial^2 \sigma_{\Delta_m}^\alpha}{\partial k_i \partial k_j} \right|_{\mathbf{k}=\mathbf{0}} = \frac{1}{\alpha^2}(\mathbf{M}_2(\alpha))_{ij}.
\end{equation}
Substituting this yields:
\begin{equation}
\langle x_i x_j \rangle(t) = \frac{\mathcal{K}_{\text{diff}}}{\alpha^2 \Gamma(\beta+1)} (\mathbf{M}_2(\alpha))_{ij}\, t^\beta.
\end{equation}
Only $E_\beta'(0)$ and the Hessian of $\sigma$ at $\mathbf{0}$ enter, so the identity is exact for all $t$, not only asymptotically.
\end{proof}

\begin{remark}[Analytical Exact Solution and Simplicial Propagation]
Fourier--Laplace solutions with Mittag-Leffler propagators are standard for continuous-time random walks and fractional diffusion \cite{metzler2000}. The contribution here is the specific compact, $S_{m-1}$-symmetric jump kernel, whose covariance $\mathbf{M}_2(\alpha)$ fixes the anisotropy of the mean squared displacement.
\end{remark}

% ==============================================================================
\section{Conclusion}
% ==============================================================================

In this paper, we have formulated and analyzed the physical dynamics governed by the Fractional Simplicial Laplacian introduced in \cite{silvafilho2026simplex}. We established the conservation laws and the modulational instability of the Nonlinear Simplicial Schr\"odinger Equation, stated the symmetry of its solitary waves as a conjecture, and derived the exact Mittag-Leffler propagator and simplicial permeability tensor for anomalous porous diffusion. Future research will explore split-step Fourier numerical solvers and experimental transport validation in lithographic microfluidic rock networks.

% Shared declarations block, \input by every chapter before its bibliography.
% Keep in sync with the "Declarations" section of master_book_unified_quantum_gravity.tex.
\section*{Declarations}

\noindent\textbf{Funding.} This research was financed in part by the Coordena\c{c}\~ao de Aperfei\c{c}oamento de Pessoal de N\'ivel Superior -- Brasil (CAPES) -- Finance Code 001.

\medskip\noindent\textbf{Competing interests.} The author declares no competing interests.

\medskip\noindent\textbf{Use of generative AI tools.} Large language model assistants (Anthropic Claude; Google Gemini, used through the Antigravity agent environment) were used extensively in preparing this work: drafting and editing text and \LaTeX{} source; proposing, expanding and revising derivations and proofs under the author's direction; writing the Python verification scripts and the \textsc{Lean 4} files; searching and checking the literature and references; and adversarial auditing of proofs, numerical claims and code. These audits found errors, some of them originally introduced with AI assistance. The errors and their corrections are recorded in the repository file \texttt{WORKPLAN.md}. AI tools are not authors. The author chose the research questions and the overall framework, reviewed the material, and is responsible for the content, including any remaining errors.

\medskip\noindent\textbf{Verification status.} Results labelled Theorem, Proposition or Lemma are proved in the text or cited as classical; statements labelled Conjecture are not proved. The numerical scripts compare formulas of the text with independent computations and are checks, not proofs. The \textsc{Lean 4} modules in \texttt{formal\_proofs\_book/Book} are a naming skeleton and do not verify the mathematics of this chapter.

\medskip\noindent\textbf{Code and data availability.} Sources, numerical scripts and the audit log are available at
\begin{center}\url{https://github.com/reinaldomsilvafilho-netizen/quantum-gravity-lean4}.\end{center}
The monograph is archived on Zenodo, \href{https://doi.org/10.5281/zenodo.22290043}{doi:10.5281/zenodo.22290043}, under the CC~BY~4.0 license.


\begin{thebibliography}{99}

\bibitem{silvafilho2026simplex} R.~M. Silva-Filho, \emph{Analytic Continuation of Pascal's Simplex: Continuous Multinomial Integrals, Polytope Boundary Recurrences, and Simplicial Fractional Calculus}, Chapter~3 of \emph{Geometry, Tensors, and Quantum Gravity}, Universidade Federal de Lavras, 2026.

\bibitem{servadei2014spectrum}
R.~Servadei and E.~Valdinoci,
\emph{On the spectrum of two different fractional operators},
Proceedings of the Royal Society of Edinburgh, Section A, \textbf{144} (2014), no.~4, 831--855,
\doi{10.1017/S0308210512001783}.

\bibitem{grisvard1985elliptic}
P.~Grisvard,
\emph{Elliptic Problems in Nonsmooth Domains},
Classics in Applied Mathematics, Vol.~69, SIAM, Philadelphia, 2011 (reprint of the 1985 edition),
\doi{10.1137/1.9781611972030}.

\bibitem{gazzola2010polyharmonic}
F.~Gazzola, H.-C. Grunau, and G.~Sweers,
\emph{Polyharmonic Boundary Value Problems},
Lecture Notes in Mathematics, Vol.~1991, Springer, Berlin, 2010,
\doi{10.1007/978-3-642-12245-3}.

\bibitem{podlubny1998}
I.~Podlubny,
\emph{Fractional Differential Equations: An Introduction to Fractional Derivatives, Fractional Differential Equations, to Methods of Their Solution and Some of Their Applications},
Mathematics in Science and Engineering, Vol.~198, Academic Press, San Diego, 1999,
\doi{10.1016/S0076-5392(99)X8001-5}.

\bibitem{samko1993}
S.~G. Samko, A.~A. Kilbas, and O.~I. Marichev,
\emph{Fractional Integrals and Derivatives: Theory and Applications},
Gordon and Breach Science Publishers, Yverdon, 1993, ISBN~978-2-88124-864-1.

\bibitem{metzler2000}
R.~Metzler and J.~Klafter,
\emph{The random walk's guide to anomalous diffusion: a fractional dynamics approach},
Physics Reports, \textbf{339} (2000), no.~1, 1--77,
\doi{10.1016/S0370-1573(00)00070-3}.

\bibitem{laskin2000}
N.~Laskin,
\emph{Fractional quantum mechanics and L\'evy path integrals},
Physics Letters A, \textbf{268} (2000), no.~4--6, 298--305,
\doi{10.1016/S0375-9601(00)00201-2}.

\end{thebibliography}

\end{document}
